%%%%%%%%%%%%%%%%%%%%%%%%%%%%%%%%%%%%%%%%%%%%%%%%%%%%%%%%%%%%
% 8.10 CATEGORICAL DATA ANALYSIS
% The Composition of Advanced Mathematics
%%%%%%%%%%%%%%%%%%%%%%%%%%%%%%%%%%%%%%%%%%%%%%%%%%%%%%%%%%%%

\section{Categorical Data Analysis}
\label{sec:categorical-data-analysis}

\prerequisite{
Probability, conditional probability, Bernoulli and binomial
distributions, descriptive statistics, sampling distributions,
confidence intervals, hypothesis testing, contingency tables,
chi-square tests, regression, and maximum likelihood estimation.
}

\learningobjectives{
By the end of this section, the reader should be able to:
\begin{itemize}
    \item distinguish categorical and quantitative variables;
    \item distinguish nominal, ordinal, and binary variables;
    \item summarize categorical variables using counts and proportions;
    \item construct and interpret contingency tables;
    \item distinguish joint, marginal, and conditional proportions;
    \item define independence for categorical variables;
    \item compute expected counts under independence;
    \item perform Pearson chi-square tests of independence;
    \item perform chi-square goodness-of-fit tests;
    \item understand likelihood-ratio chi-square statistics;
    \item compute and interpret risk, odds, relative risk, and odds ratios;
    \item construct large-sample intervals for odds ratios;
    \item understand Fisher's exact test conceptually;
    \item understand McNemar's test for paired binary data;
    \item understand logistic regression;
    \item interpret logistic regression coefficients through odds ratios;
    \item estimate predicted probabilities from logistic models;
    \item understand multiple logistic regression;
    \item use indicator variables for categorical predictors;
    \item understand interaction terms in logistic regression;
    \item understand likelihood-ratio, Wald, and score tests in logistic
          models;
    \item understand deviance and model comparison;
    \item understand goodness-of-fit diagnostics;
    \item understand multinomial logistic regression conceptually;
    \item understand ordinal logistic regression conceptually;
    \item understand proportional-odds models;
    \item understand log-linear models for contingency tables;
    \item distinguish prospective and retrospective sampling;
    \item understand stratified categorical analysis;
    \item understand the Mantel--Haenszel estimator conceptually;
    \item recognize Simpson's paradox;
    \item distinguish confounding from interaction;
    \item understand exact and asymptotic categorical inference;
    \item connect categorical-data methods with generalized linear
          models, regression, experimental design, and statistical
          learning.
\end{itemize}
}

%%%%%%%%%%%%%%%%%%%%%%%%%%%%%%%%%%%%%%%%%%%%%%%%%%%%%%%%%%%%
\subsection{Categorical Variables}
\label{subsec:categorical-variables}
%%%%%%%%%%%%%%%%%%%%%%%%%%%%%%%%%%%%%%%%%%%%%%%%%%%%%%%%%%%%

A categorical variable records membership in one of several categories.

Examples include

\[
\boxed{
\text{yes/no},
}
\]

\[
\boxed{
\text{blood type},
}
\]

\[
\boxed{
\text{political region},
}
\]

or

\[
\boxed{
\text{disease stage}.
}
\]

The mathematical analysis focuses primarily on counts and probabilities
rather than distances between numerical values.

%%%%%%%%%%%%%%%%%%%%%%%%%%%%%%%%%%%%%%%%%%%%%%%%%%%%%%%%%%%%
\subsection{Nominal Variables}
\label{subsec:nominal-variables}
%%%%%%%%%%%%%%%%%%%%%%%%%%%%%%%%%%%%%%%%%%%%%%%%%%%%%%%%%%%%

A nominal variable has categories with no intrinsic ordering.

Examples include

\[
\boxed{
\text{blood type}
}
\]

or

\[
\boxed{
\text{transportation mode}.
}
\]

Numerical labels may be assigned for computation, but those labels do not
imply quantitative spacing.

%%%%%%%%%%%%%%%%%%%%%%%%%%%%%%%%%%%%%%%%%%%%%%%%%%%%%%%%%%%%
\subsection{Ordinal Variables}
\label{subsec:ordinal-variables}
%%%%%%%%%%%%%%%%%%%%%%%%%%%%%%%%%%%%%%%%%%%%%%%%%%%%%%%%%%%%

An ordinal variable has categories with a natural ordering.

Examples include

\[
\boxed{
\text{low},
\quad
\text{medium},
\quad
\text{high},
}
\]

or

\[
\boxed{
\text{strongly disagree}
<
\text{disagree}
<
\text{neutral}
<
\text{agree}
<
\text{strongly agree}.
}
\]

The order is meaningful, but distances between categories need not be
equal.

%%%%%%%%%%%%%%%%%%%%%%%%%%%%%%%%%%%%%%%%%%%%%%%%%%%%%%%%%%%%
\subsection{Binary Variables}
\label{subsec:binary-variables}
%%%%%%%%%%%%%%%%%%%%%%%%%%%%%%%%%%%%%%%%%%%%%%%%%%%%%%%%%%%%

A binary variable has two categories.

It is commonly coded as

\[
\boxed{
Y
=
\begin{cases}
1,&\text{success},\\
0,&\text{failure}.
\end{cases}
}
\]

Then

\[
Y\sim\operatorname{Bernoulli}(p)
\]

under a Bernoulli model.

%%%%%%%%%%%%%%%%%%%%%%%%%%%%%%%%%%%%%%%%%%%%%%%%%%%%%%%%%%%%
\subsection{Frequency Tables}
\label{subsec:categorical-frequency-tables}
%%%%%%%%%%%%%%%%%%%%%%%%%%%%%%%%%%%%%%%%%%%%%%%%%%%%%%%%%%%%

Suppose a categorical variable has categories

\[
1,\ldots,k.
\]

Let

\[
n_i
\]

be the number of observations in category \(i\).

Then

\[
\boxed{
\sum_{i=1}^{k}n_i=n.
}
\]

The sample proportion in category \(i\) is

\[
\boxed{
\widehat{p}_i
=
\frac{n_i}{n}.
}
\]

%%%%%%%%%%%%%%%%%%%%%%%%%%%%%%%%%%%%%%%%%%%%%%%%%%%%%%%%%%%%
\subsection{Contingency Tables}
\label{subsec:contingency-tables}
%%%%%%%%%%%%%%%%%%%%%%%%%%%%%%%%%%%%%%%%%%%%%%%%%%%%%%%%%%%%

A contingency table summarizes the joint distribution of two
categorical variables.

Suppose variable \(A\) has \(r\) categories and variable \(B\) has \(c\)
categories.

Let

\[
n_{ij}
\]

denote the count in row \(i\), column \(j\).

Then

\[
\boxed{
n
=
\sum_{i=1}^{r}
\sum_{j=1}^{c}
n_{ij}.
}
\]

%%%%%%%%%%%%%%%%%%%%%%%%%%%%%%%%%%%%%%%%%%%%%%%%%%%%%%%%%%%%
\subsection{Marginal Counts}
\label{subsec:marginal-counts}
%%%%%%%%%%%%%%%%%%%%%%%%%%%%%%%%%%%%%%%%%%%%%%%%%%%%%%%%%%%%

The row total is

\[
\boxed{
n_{i\cdot}
=
\sum_{j=1}^{c}n_{ij}.
}
\]

The column total is

\[
\boxed{
n_{\cdot j}
=
\sum_{i=1}^{r}n_{ij}.
}
\]

These are called marginal counts.

%%%%%%%%%%%%%%%%%%%%%%%%%%%%%%%%%%%%%%%%%%%%%%%%%%%%%%%%%%%%
\subsection{Joint Proportions}
\label{subsec:joint-proportions-categorical}
%%%%%%%%%%%%%%%%%%%%%%%%%%%%%%%%%%%%%%%%%%%%%%%%%%%%%%%%%%%%

The joint sample proportion for cell \((i,j)\) is

\[
\boxed{
\widehat{p}_{ij}
=
\frac{n_{ij}}{n}.
}
\]

The marginal proportions are

\[
\boxed{
\widehat{p}_{i\cdot}
=
\frac{n_{i\cdot}}{n},
}
\]

and

\[
\boxed{
\widehat{p}_{\cdot j}
=
\frac{n_{\cdot j}}{n}.
}
\]

%%%%%%%%%%%%%%%%%%%%%%%%%%%%%%%%%%%%%%%%%%%%%%%%%%%%%%%%%%%%
\subsection{Conditional Proportions}
\label{subsec:conditional-proportions-categorical}
%%%%%%%%%%%%%%%%%%%%%%%%%%%%%%%%%%%%%%%%%%%%%%%%%%%%%%%%%%%%

The conditional proportion of column \(j\) within row \(i\) is

\[
\boxed{
\frac{n_{ij}}{n_{i\cdot}}.
}
\]

Likewise, the conditional proportion of row \(i\) within column \(j\) is

\[
\boxed{
\frac{n_{ij}}{n_{\cdot j}}.
}
\]

Conditional proportions are often more informative than raw counts when
group sizes differ.

%%%%%%%%%%%%%%%%%%%%%%%%%%%%%%%%%%%%%%%%%%%%%%%%%%%%%%%%%%%%
\subsection{Independence of Categorical Variables}
\label{subsec:categorical-independence}
%%%%%%%%%%%%%%%%%%%%%%%%%%%%%%%%%%%%%%%%%%%%%%%%%%%%%%%%%%%%

Two categorical variables \(A\) and \(B\) are independent if

\[
\boxed{
P(A=i,B=j)
=
P(A=i)P(B=j)
}
\]

for every pair of categories.

Equivalently,

\[
\boxed{
P(B=j\mid A=i)
=
P(B=j)
}
\]

whenever the conditional probability is defined.

%%%%%%%%%%%%%%%%%%%%%%%%%%%%%%%%%%%%%%%%%%%%%%%%%%%%%%%%%%%%
\subsection{Expected Counts Under Independence}
\label{subsec:expected-counts-categorical}
%%%%%%%%%%%%%%%%%%%%%%%%%%%%%%%%%%%%%%%%%%%%%%%%%%%%%%%%%%%%

Under independence, the estimated expected count in cell \((i,j)\) is

\[
\boxed{
E_{ij}
=
\frac{
n_{i\cdot}n_{\cdot j}
}{
n
}.
}
\]

This is obtained by multiplying estimated marginal probabilities:

\[
\frac{n_{i\cdot}}{n}
\frac{n_{\cdot j}}{n}
\]

and multiplying by \(n\).

%%%%%%%%%%%%%%%%%%%%%%%%%%%%%%%%%%%%%%%%%%%%%%%%%%%%%%%%%%%%
\subsection{Worked Example: Expected Cell Count}
\label{subsec:worked-expected-cell-count-categorical}
%%%%%%%%%%%%%%%%%%%%%%%%%%%%%%%%%%%%%%%%%%%%%%%%%%%%%%%%%%%%

\begin{workedexamplebox}{Expected Count Under Independence}

Suppose a contingency table has grand total

\[
n=300.
\]

A row total is

\[
120
\]

and a column total is

\[
75.
\]

Then the expected count is

\[
E
=
\frac{
120(75)
}{
300
}.
\]

Therefore,

\begin{answerbox}
\[
\boxed{
E=30.
}
\]
\end{answerbox}

\end{workedexamplebox}

%%%%%%%%%%%%%%%%%%%%%%%%%%%%%%%%%%%%%%%%%%%%%%%%%%%%%%%%%%%%
\subsection{Pearson Chi-Square Test of Independence}
\label{subsec:pearson-chi-square-independence}
%%%%%%%%%%%%%%%%%%%%%%%%%%%%%%%%%%%%%%%%%%%%%%%%%%%%%%%%%%%%

The Pearson statistic is

\[
\boxed{
X^2
=
\sum_{i=1}^{r}
\sum_{j=1}^{c}
\frac{
(O_{ij}-E_{ij})^2
}{
E_{ij}
}.
}
\]

Under the null hypothesis of independence and suitable large-sample
conditions,

\[
\boxed{
X^2
\approx
\chi^2_{(r-1)(c-1)}.
}
\]

%%%%%%%%%%%%%%%%%%%%%%%%%%%%%%%%%%%%%%%%%%%%%%%%%%%%%%%%%%%%
\subsection{Hypotheses for Independence}
\label{subsec:hypotheses-independence-categorical}
%%%%%%%%%%%%%%%%%%%%%%%%%%%%%%%%%%%%%%%%%%%%%%%%%%%%%%%%%%%%

The null hypothesis is

\[
\boxed{
H_0:
A
\text{ and }
B
\text{ are independent}.
}
\]

The alternative is

\[
\boxed{
H_A:
A
\text{ and }
B
\text{ are associated}.
}
\]

A significant result does not describe the direction or strength of the
association by itself.

%%%%%%%%%%%%%%%%%%%%%%%%%%%%%%%%%%%%%%%%%%%%%%%%%%%%%%%%%%%%
\subsection{Pearson Residuals}
\label{subsec:pearson-residuals-categorical}
%%%%%%%%%%%%%%%%%%%%%%%%%%%%%%%%%%%%%%%%%%%%%%%%%%%%%%%%%%%%

A Pearson residual is

\[
\boxed{
r_{ij}
=
\frac{
O_{ij}-E_{ij}
}{
\sqrt{E_{ij}}
}.
}
\]

Cells with large absolute residuals contribute strongly to the overall
chi-square statistic.

Indeed,

\[
\boxed{
X^2
=
\sum_{i,j}r_{ij}^2.
}
\]

%%%%%%%%%%%%%%%%%%%%%%%%%%%%%%%%%%%%%%%%%%%%%%%%%%%%%%%%%%%%
\subsection{Goodness-of-Fit for One Categorical Variable}
\label{subsec:categorical-goodness-fit}
%%%%%%%%%%%%%%%%%%%%%%%%%%%%%%%%%%%%%%%%%%%%%%%%%%%%%%%%%%%%

Suppose category probabilities under the null are

\[
p_1,\ldots,p_k.
\]

Then

\[
\boxed{
E_i
=
np_i.
}
\]

The Pearson statistic is

\[
\boxed{
X^2
=
\sum_{i=1}^{k}
\frac{
(O_i-E_i)^2
}{
E_i
}.
}
\]

If the probabilities are completely specified,

\[
\boxed{
df=k-1.
}
\]

%%%%%%%%%%%%%%%%%%%%%%%%%%%%%%%%%%%%%%%%%%%%%%%%%%%%%%%%%%%%
\subsection{Likelihood-Ratio Chi-Square Statistic}
\label{subsec:likelihood-ratio-chi-square}
%%%%%%%%%%%%%%%%%%%%%%%%%%%%%%%%%%%%%%%%%%%%%%%%%%%%%%%%%%%%

Another categorical-data statistic is the likelihood-ratio statistic

\[
\boxed{
G^2
=
2
\sum
O_i
\log
\left(
\frac{
O_i
}{
E_i
}
\right).
}
\]

For contingency tables,

\[
\boxed{
G^2
=
2
\sum_{i,j}
O_{ij}
\log
\left(
\frac{
O_{ij}
}{
E_{ij}
}
\right).
}
\]

Cells with \(O_{ij}=0\) contribute \(0\) by continuity.

%%%%%%%%%%%%%%%%%%%%%%%%%%%%%%%%%%%%%%%%%%%%%%%%%%%%%%%%%%%%
\subsection{Pearson \(X^2\) Versus \(G^2\)}
\label{subsec:pearson-versus-likelihood-ratio}
%%%%%%%%%%%%%%%%%%%%%%%%%%%%%%%%%%%%%%%%%%%%%%%%%%%%%%%%%%%%

For large samples under regular conditions,

\[
\boxed{
X^2
}
\]

and

\[
\boxed{
G^2
}
\]

have the same asymptotic chi-square distribution.

They arise from different discrepancy measures but are asymptotically
equivalent under the null.

%%%%%%%%%%%%%%%%%%%%%%%%%%%%%%%%%%%%%%%%%%%%%%%%%%%%%%%%%%%%
\subsection{The \(2\times2\) Table}
\label{subsec:two-by-two-table}
%%%%%%%%%%%%%%%%%%%%%%%%%%%%%%%%%%%%%%%%%%%%%%%%%%%%%%%%%%%%

A particularly important contingency table is

\[
\begin{array}{c|cc|c}
&
Y=1
&
Y=0
&
\text{Total}
\\
\hline
X=1
&
a
&
b
&
a+b
\\
X=0
&
c
&
d
&
c+d
\\
\hline
\text{Total}
&
a+c
&
b+d
&
n
\end{array}
\]

This structure supports several measures of association.

%%%%%%%%%%%%%%%%%%%%%%%%%%%%%%%%%%%%%%%%%%%%%%%%%%%%%%%%%%%%
\subsection{Risk}
\label{subsec:risk-categorical}
%%%%%%%%%%%%%%%%%%%%%%%%%%%%%%%%%%%%%%%%%%%%%%%%%%%%%%%%%%%%

For the exposed or treatment group \(X=1\),

\[
\boxed{
\widehat{p}_1
=
\frac{
a
}{
a+b
}.
}
\]

For the unexposed or comparison group \(X=0\),

\[
\boxed{
\widehat{p}_0
=
\frac{
c
}{
c+d
}.
}
\]

These are estimated risks or event probabilities.

%%%%%%%%%%%%%%%%%%%%%%%%%%%%%%%%%%%%%%%%%%%%%%%%%%%%%%%%%%%%
\subsection{Risk Difference}
\label{subsec:risk-difference}
%%%%%%%%%%%%%%%%%%%%%%%%%%%%%%%%%%%%%%%%%%%%%%%%%%%%%%%%%%%%

The risk difference is

\[
\boxed{
RD
=
p_1-p_0.
}
\]

The sample estimate is

\[
\boxed{
\widehat{RD}
=
\widehat{p}_1-\widehat{p}_0.
}
\]

A value of

\[
0
\]

corresponds to equal risks.

%%%%%%%%%%%%%%%%%%%%%%%%%%%%%%%%%%%%%%%%%%%%%%%%%%%%%%%%%%%%
\subsection{Relative Risk}
\label{subsec:relative-risk}
%%%%%%%%%%%%%%%%%%%%%%%%%%%%%%%%%%%%%%%%%%%%%%%%%%%%%%%%%%%%

The relative risk is

\[
\boxed{
RR
=
\frac{
p_1
}{
p_0
}.
}
\]

Its estimate is

\[
\boxed{
\widehat{RR}
=
\frac{
a/(a+b)
}{
c/(c+d)
}.
}
\]

A value of

\[
\boxed{
RR=1
}
\]

indicates equal risks.

%%%%%%%%%%%%%%%%%%%%%%%%%%%%%%%%%%%%%%%%%%%%%%%%%%%%%%%%%%%%
\subsection{Odds}
\label{subsec:odds}
%%%%%%%%%%%%%%%%%%%%%%%%%%%%%%%%%%%%%%%%%%%%%%%%%%%%%%%%%%%%

For probability \(p\), the odds are

\[
\boxed{
\operatorname{odds}
=
\frac{
p
}{
1-p
}.
}
\]

If

\[
p=\frac12,
\]

then the odds are

\[
1.
\]

If

\[
p=0.8,
\]

then

\[
\boxed{
\operatorname{odds}
=
\frac{0.8}{0.2}
=
4.
}
\]

%%%%%%%%%%%%%%%%%%%%%%%%%%%%%%%%%%%%%%%%%%%%%%%%%%%%%%%%%%%%
\subsection{Odds Ratio}
\label{subsec:odds-ratio}
%%%%%%%%%%%%%%%%%%%%%%%%%%%%%%%%%%%%%%%%%%%%%%%%%%%%%%%%%%%%

The odds ratio compares odds between two groups:

\[
\boxed{
OR
=
\frac{
p_1/(1-p_1)
}{
p_0/(1-p_0)
}.
}
\]

For a \(2\times2\) table,

\[
\boxed{
\widehat{OR}
=
\frac{
ad
}{
bc
},
}
\]

provided the relevant cell counts are nonzero.

%%%%%%%%%%%%%%%%%%%%%%%%%%%%%%%%%%%%%%%%%%%%%%%%%%%%%%%%%%%%
\subsection{Worked Example: Odds Ratio}
\label{subsec:worked-odds-ratio}
%%%%%%%%%%%%%%%%%%%%%%%%%%%%%%%%%%%%%%%%%%%%%%%%%%%%%%%%%%%%

\begin{workedexamplebox}{Computing an Odds Ratio}

Suppose

\[
a=40,
\qquad
b=20,
\qquad
c=30,
\qquad
d=60.
\]

Then

\[
\widehat{OR}
=
\frac{
40(60)
}{
20(30)
}.
\]

Therefore,

\begin{answerbox}
\[
\boxed{
\widehat{OR}=4.
}
\]
\end{answerbox}

The estimated odds of the event are four times as large in the first
group.

\end{workedexamplebox}

%%%%%%%%%%%%%%%%%%%%%%%%%%%%%%%%%%%%%%%%%%%%%%%%%%%%%%%%%%%%
\subsection{Risk Ratio Versus Odds Ratio}
\label{subsec:risk-ratio-versus-odds-ratio}
%%%%%%%%%%%%%%%%%%%%%%%%%%%%%%%%%%%%%%%%%%%%%%%%%%%%%%%%%%%%

Relative risk compares probabilities directly:

\[
RR
=
\frac{p_1}{p_0}.
\]

The odds ratio compares odds:

\[
OR
=
\frac{
p_1/(1-p_1)
}{
p_0/(1-p_0)
}.
\]

They are not generally equal.

When events are rare,

\[
\boxed{
OR
\approx
RR.
}
\]

For common events, the difference can be substantial.

%%%%%%%%%%%%%%%%%%%%%%%%%%%%%%%%%%%%%%%%%%%%%%%%%%%%%%%%%%%%
\subsection{Log Odds Ratio}
\label{subsec:log-odds-ratio}
%%%%%%%%%%%%%%%%%%%%%%%%%%%%%%%%%%%%%%%%%%%%%%%%%%%%%%%%%%%%

The logarithm of the odds ratio is

\[
\boxed{
\log OR.
}
\]

Its null value is

\[
\boxed{
0,
}
\]

because

\[
\log 1=0.
\]

For a \(2\times2\) table,

\[
\boxed{
\log\widehat{OR}
=
\log a
+
\log d
-
\log b
-
\log c.
}
\]

%%%%%%%%%%%%%%%%%%%%%%%%%%%%%%%%%%%%%%%%%%%%%%%%%%%%%%%%%%%%
\subsection{Approximate Standard Error of the Log Odds Ratio}
\label{subsec:se-log-odds-ratio}
%%%%%%%%%%%%%%%%%%%%%%%%%%%%%%%%%%%%%%%%%%%%%%%%%%%%%%%%%%%%

For independent counts under standard large-sample conditions,

\[
\boxed{
SE
(
\log\widehat{OR}
)
\approx
\sqrt{
\frac1a
+
\frac1b
+
\frac1c
+
\frac1d
}.
}
\]

Therefore an approximate confidence interval on the log scale is

\[
\boxed{
\log\widehat{OR}
\pm
z^*
SE
(
\log\widehat{OR}
).
}
\]

Exponentiating the endpoints gives an interval for \(OR\).

%%%%%%%%%%%%%%%%%%%%%%%%%%%%%%%%%%%%%%%%%%%%%%%%%%%%%%%%%%%%
\subsection{Worked Example: Odds-Ratio Confidence Interval Form}
\label{subsec:worked-or-ci}
%%%%%%%%%%%%%%%%%%%%%%%%%%%%%%%%%%%%%%%%%%%%%%%%%%%%%%%%%%%%

\begin{workedexamplebox}{Large-Sample Odds-Ratio Interval}

Suppose

\[
a=50,
\quad
b=25,
\quad
c=20,
\quad
d=40.
\]

Then

\[
\widehat{OR}
=
\frac{
50(40)
}{
25(20)
}
=
4.
\]

The estimated log-scale standard error is

\[
SE
=
\sqrt{
\frac1{50}
+
\frac1{25}
+
\frac1{20}
+
\frac1{40}
}.
\]

A \(95\%\) interval is first constructed as

\[
\log4
\pm
1.96SE.
\]

Then exponentiate.

\begin{answerbox}
\[
\boxed{
CI_{OR}
=
\left(
e^{\log4-1.96SE},
e^{\log4+1.96SE}
\right).
}
\]
\end{answerbox}

\end{workedexamplebox}

%%%%%%%%%%%%%%%%%%%%%%%%%%%%%%%%%%%%%%%%%%%%%%%%%%%%%%%%%%%%
\subsection{Case-Control Sampling and Odds Ratios}
\label{subsec:case-control-odds-ratio}
%%%%%%%%%%%%%%%%%%%%%%%%%%%%%%%%%%%%%%%%%%%%%%%%%%%%%%%%%%%%

In a case-control study, observations may be sampled according to outcome
status rather than exposure status.

Then absolute risks generally cannot be estimated directly from the
case-control sample without additional information.

However, the exposure odds ratio can still estimate the corresponding
population odds ratio under the standard sampling framework.

This is one reason odds ratios are central in epidemiology.

%%%%%%%%%%%%%%%%%%%%%%%%%%%%%%%%%%%%%%%%%%%%%%%%%%%%%%%%%%%%
\subsection{Prospective and Retrospective Sampling}
\label{subsec:prospective-retrospective}
%%%%%%%%%%%%%%%%%%%%%%%%%%%%%%%%%%%%%%%%%%%%%%%%%%%%%%%%%%%%

Prospective sampling often selects units based on explanatory variables
or treatment groups and then observes outcomes.

Retrospective sampling may select units based on outcomes and examine
past exposures.

The sampling scheme determines which probabilities are directly
estimable.

%%%%%%%%%%%%%%%%%%%%%%%%%%%%%%%%%%%%%%%%%%%%%%%%%%%%%%%%%%%%
\subsection{Fisher's Exact Test}
\label{subsec:fishers-exact-test}
%%%%%%%%%%%%%%%%%%%%%%%%%%%%%%%%%%%%%%%%%%%%%%%%%%%%%%%%%%%%

For a \(2\times2\) table with fixed margins, the conditional distribution
of one cell under independence is hypergeometric.

This leads to Fisher's exact test.

If the margins are fixed, the probability of observing cell count \(A=a\)
is

\[
\boxed{
P(A=a)
=
\frac{
\binom{a+b}{a}
\binom{c+d}{c}
}{
\binom{n}{a+c}
},
}
\]

with equivalent forms depending on which margins are represented.

%%%%%%%%%%%%%%%%%%%%%%%%%%%%%%%%%%%%%%%%%%%%%%%%%%%%%%%%%%%%
\subsection{When Fisher's Exact Test Is Useful}
\label{subsec:when-fisher-exact}
%%%%%%%%%%%%%%%%%%%%%%%%%%%%%%%%%%%%%%%%%%%%%%%%%%%%%%%%%%%%

Fisher's exact test is especially useful when sample sizes are small or
expected counts are too small for a chi-square approximation.

It provides exact conditional inference under the specified fixed-margin
framework.

Two-sided exact \(p\)-values require a rule for defining tables at least
as extreme as the observed table.

%%%%%%%%%%%%%%%%%%%%%%%%%%%%%%%%%%%%%%%%%%%%%%%%%%%%%%%%%%%%
\subsection{Matched Binary Data}
\label{subsec:matched-binary-data}
%%%%%%%%%%%%%%%%%%%%%%%%%%%%%%%%%%%%%%%%%%%%%%%%%%%%%%%%%%%%

Suppose each subject is measured twice or matched into pairs.

A paired \(2\times2\) table may be written

\[
\begin{array}{c|cc}
&
\text{After }1
&
\text{After }0
\\
\hline
\text{Before }1
&
a
&
b
\\
\text{Before }0
&
c
&
d
\end{array}
\]

The concordant pairs are

\[
a
\]

and

\[
d.
\]

The discordant pairs are

\[
b
\]

and

\[
c.
\]

%%%%%%%%%%%%%%%%%%%%%%%%%%%%%%%%%%%%%%%%%%%%%%%%%%%%%%%%%%%%
\subsection{McNemar's Test}
\label{subsec:mcnemar-test}
%%%%%%%%%%%%%%%%%%%%%%%%%%%%%%%%%%%%%%%%%%%%%%%%%%%%%%%%%%%%

McNemar's test evaluates whether the probabilities of the two types of
discordant pairs are equal.

The large-sample statistic is

\[
\boxed{
X^2
=
\frac{
(b-c)^2
}{
b+c
}.
}
\]

Under the null and suitable conditions,

\[
\boxed{
X^2
\approx
\chi^2_1.
}
\]

An exact binomial version can be used for small numbers of discordant
pairs.

%%%%%%%%%%%%%%%%%%%%%%%%%%%%%%%%%%%%%%%%%%%%%%%%%%%%%%%%%%%%
\subsection{Why Only Discordant Pairs Matter}
\label{subsec:discordant-pairs}
%%%%%%%%%%%%%%%%%%%%%%%%%%%%%%%%%%%%%%%%%%%%%%%%%%%%%%%%%%%%

Concordant pairs do not indicate a change in binary status.

Only pairs that change from

\[
1\to0
\]

or

\[
0\to1
\]

contain information about directional change.

Thus McNemar's test is based on

\[
\boxed{
b
\text{ and }
c.
}
\]

%%%%%%%%%%%%%%%%%%%%%%%%%%%%%%%%%%%%%%%%%%%%%%%%%%%%%%%%%%%%
\subsection{Binary Regression}
\label{subsec:binary-regression}
%%%%%%%%%%%%%%%%%%%%%%%%%%%%%%%%%%%%%%%%%%%%%%%%%%%%%%%%%%%%

Suppose

\[
Y\in\{0,1\}.
\]

Linear regression is often inappropriate because it may predict values
outside

\[
[0,1].
\]

Instead, model the probability

\[
\boxed{
p(x)
=
P(Y=1\mid X=x).
}
\]

Logistic regression is the most common model.

%%%%%%%%%%%%%%%%%%%%%%%%%%%%%%%%%%%%%%%%%%%%%%%%%%%%%%%%%%%%
\subsection{Logit Function}
\label{subsec:logit-function}
%%%%%%%%%%%%%%%%%%%%%%%%%%%%%%%%%%%%%%%%%%%%%%%%%%%%%%%%%%%%

The logit function is

\[
\boxed{
\operatorname{logit}(p)
=
\log
\left(
\frac{
p
}{
1-p
}
\right).
}
\]

It maps

\[
(0,1)
\]

onto

\[
\mathbb{R}.
\]

Thus a linear predictor can model the transformed probability.

%%%%%%%%%%%%%%%%%%%%%%%%%%%%%%%%%%%%%%%%%%%%%%%%%%%%%%%%%%%%
\subsection{Simple Logistic Regression}
\label{subsec:simple-logistic-regression}
%%%%%%%%%%%%%%%%%%%%%%%%%%%%%%%%%%%%%%%%%%%%%%%%%%%%%%%%%%%%

Simple logistic regression assumes

\[
\boxed{
\log
\left(
\frac{
p(x)
}{
1-p(x)
}
\right)
=
\beta_0+\beta_1x.
}
\]

Solving for \(p(x)\),

\[
\boxed{
p(x)
=
\frac{
e^{\beta_0+\beta_1x}
}{
1+
e^{\beta_0+\beta_1x}
}.
}
\]

This guarantees

\[
0<p(x)<1.
\]

%%%%%%%%%%%%%%%%%%%%%%%%%%%%%%%%%%%%%%%%%%%%%%%%%%%%%%%%%%%%
\subsection{Logistic Curve}
\label{subsec:logistic-curve}
%%%%%%%%%%%%%%%%%%%%%%%%%%%%%%%%%%%%%%%%%%%%%%%%%%%%%%%%%%%%

The function

\[
p(x)
=
\frac{
e^{\eta}
}{
1+e^{\eta}
},
\qquad
\eta=\beta_0+\beta_1x,
\]

has an S-shaped form.

If

\[
\beta_1>0,
\]

the probability increases with \(x\).

If

\[
\beta_1<0,
\]

the probability decreases with \(x\).

%%%%%%%%%%%%%%%%%%%%%%%%%%%%%%%%%%%%%%%%%%%%%%%%%%%%%%%%%%%%
\subsection{Interpretation of Logistic Regression Slope}
\label{subsec:logistic-slope-interpretation}
%%%%%%%%%%%%%%%%%%%%%%%%%%%%%%%%%%%%%%%%%%%%%%%%%%%%%%%%%%%%

From

\[
\operatorname{logit}(p)
=
\beta_0+\beta_1x,
\]

a one-unit increase in \(x\) increases the log-odds by

\[
\boxed{
\beta_1.
}
\]

Exponentiating,

\[
\boxed{
e^{\beta_1}
}
\]

is the odds ratio associated with a one-unit increase in \(x\).

%%%%%%%%%%%%%%%%%%%%%%%%%%%%%%%%%%%%%%%%%%%%%%%%%%%%%%%%%%%%
\subsection{Worked Example: Logistic Odds Ratio}
\label{subsec:worked-logistic-odds-ratio}
%%%%%%%%%%%%%%%%%%%%%%%%%%%%%%%%%%%%%%%%%%%%%%%%%%%%%%%%%%%%

\begin{workedexamplebox}{Interpreting a Logistic Coefficient}

Suppose a logistic regression coefficient is

\[
\beta_1=\log 2.
\]

Then

\[
e^{\beta_1}=2.
\]

\begin{answerbox}
\[
\boxed{
\text{A one-unit increase in }X
\text{ doubles the modeled odds of }Y=1.
}
\]
\end{answerbox}

This does not mean the probability itself doubles.

\end{workedexamplebox}

%%%%%%%%%%%%%%%%%%%%%%%%%%%%%%%%%%%%%%%%%%%%%%%%%%%%%%%%%%%%
\subsection{Intercept in Logistic Regression}
\label{subsec:logistic-intercept}
%%%%%%%%%%%%%%%%%%%%%%%%%%%%%%%%%%%%%%%%%%%%%%%%%%%%%%%%%%%%

At

\[
x=0,
\]

the log-odds are

\[
\boxed{
\beta_0.
}
\]

Thus the odds are

\[
\boxed{
e^{\beta_0},
}
\]

and the corresponding probability is

\[
\boxed{
\frac{
e^{\beta_0}
}{
1+e^{\beta_0}
}.
}
\]

As in linear regression, the intercept may not be scientifically
meaningful if \(x=0\) is outside the relevant range.

%%%%%%%%%%%%%%%%%%%%%%%%%%%%%%%%%%%%%%%%%%%%%%%%%%%%%%%%%%%%
\subsection{Worked Example: Predicted Probability}
\label{subsec:worked-logistic-probability}
%%%%%%%%%%%%%%%%%%%%%%%%%%%%%%%%%%%%%%%%%%%%%%%%%%%%%%%%%%%%

\begin{workedexamplebox}{Converting Log-Odds to Probability}

Suppose

\[
\operatorname{logit}(p)
=
-2+0.5x.
\]

At

\[
x=4,
\]

the linear predictor is

\[
-2+0.5(4)=0.
\]

Therefore,

\[
p
=
\frac{
e^0
}{
1+e^0
}
=
\frac12.
\]

\begin{answerbox}
\[
\boxed{
p=0.5.
}
\]
\end{answerbox}

\end{workedexamplebox}

%%%%%%%%%%%%%%%%%%%%%%%%%%%%%%%%%%%%%%%%%%%%%%%%%%%%%%%%%%%%
\subsection{Likelihood for Logistic Regression}
\label{subsec:logistic-likelihood}
%%%%%%%%%%%%%%%%%%%%%%%%%%%%%%%%%%%%%%%%%%%%%%%%%%%%%%%%%%%%

For independent binary responses,

\[
Y_i
\sim
\operatorname{Bernoulli}(p_i),
\]

where

\[
p_i
=
\frac{
e^{\mathbf{x}_i^T\boldsymbol{\beta}}
}{
1+
e^{\mathbf{x}_i^T\boldsymbol{\beta}}
}.
\]

The likelihood is

\[
\boxed{
L(\boldsymbol{\beta})
=
\prod_{i=1}^{n}
p_i^{y_i}
(1-p_i)^{1-y_i}.
}
\]

%%%%%%%%%%%%%%%%%%%%%%%%%%%%%%%%%%%%%%%%%%%%%%%%%%%%%%%%%%%%
\subsection{Log-Likelihood for Logistic Regression}
\label{subsec:logistic-log-likelihood}
%%%%%%%%%%%%%%%%%%%%%%%%%%%%%%%%%%%%%%%%%%%%%%%%%%%%%%%%%%%%

Taking logarithms,

\[
\boxed{
\ell(\boldsymbol{\beta})
=
\sum_{i=1}^{n}
\left[
y_i\log p_i
+
(1-y_i)\log(1-p_i)
\right].
}
\]

Unlike ordinary least squares, the coefficient estimates usually do not
have a simple closed-form expression.

Numerical optimization is used.

%%%%%%%%%%%%%%%%%%%%%%%%%%%%%%%%%%%%%%%%%%%%%%%%%%%%%%%%%%%%
\subsection{Multiple Logistic Regression}
\label{subsec:multiple-logistic-regression}
%%%%%%%%%%%%%%%%%%%%%%%%%%%%%%%%%%%%%%%%%%%%%%%%%%%%%%%%%%%%

With predictors

\[
X_1,\ldots,X_p,
\]

the model is

\[
\boxed{
\operatorname{logit}(p)
=
\beta_0
+
\beta_1X_1
+
\cdots
+
\beta_pX_p.
}
\]

For predictor \(X_j\),

\[
\boxed{
e^{\beta_j}
}
\]

is the multiplicative change in odds associated with a one-unit increase
in \(X_j\), holding the other modeled predictors fixed.

%%%%%%%%%%%%%%%%%%%%%%%%%%%%%%%%%%%%%%%%%%%%%%%%%%%%%%%%%%%%
\subsection{Categorical Predictors in Logistic Regression}
\label{subsec:categorical-predictors-logistic}
%%%%%%%%%%%%%%%%%%%%%%%%%%%%%%%%%%%%%%%%%%%%%%%%%%%%%%%%%%%%

Categorical predictors can be included using indicator variables.

Suppose group \(B\) is compared with reference group \(A\):

\[
D
=
\begin{cases}
1,&B,\\
0,&A.
\end{cases}
\]

Then

\[
\boxed{
\operatorname{logit}(p)
=
\beta_0+\beta_1D.
}
\]

The group odds ratio is

\[
\boxed{
e^{\beta_1}.
}
\]

%%%%%%%%%%%%%%%%%%%%%%%%%%%%%%%%%%%%%%%%%%%%%%%%%%%%%%%%%%%%
\subsection{Interaction in Logistic Regression}
\label{subsec:logistic-interaction}
%%%%%%%%%%%%%%%%%%%%%%%%%%%%%%%%%%%%%%%%%%%%%%%%%%%%%%%%%%%%

Consider

\[
\boxed{
\operatorname{logit}(p)
=
\beta_0
+
\beta_1X
+
\beta_2Z
+
\beta_3XZ.
}
\]

The log-odds effect of \(X\) is

\[
\boxed{
\beta_1+\beta_3Z.
}
\]

Thus the odds ratio for a one-unit increase in \(X\) depends on \(Z\):

\[
\boxed{
\exp
(
\beta_1+\beta_3Z
).
}
\]

%%%%%%%%%%%%%%%%%%%%%%%%%%%%%%%%%%%%%%%%%%%%%%%%%%%%%%%%%%%%
\subsection{Wald Test in Logistic Regression}
\label{subsec:wald-test-logistic}
%%%%%%%%%%%%%%%%%%%%%%%%%%%%%%%%%%%%%%%%%%%%%%%%%%%%%%%%%%%%

A coefficient can be tested using

\[
\boxed{
Z
=
\frac{
\widehat{\beta}_j
}{
SE(\widehat{\beta}_j)
}
}
\]

for

\[
H_0:\beta_j=0.
\]

Under large-sample regularity conditions,

\[
Z
\]

is approximately standard normal.

%%%%%%%%%%%%%%%%%%%%%%%%%%%%%%%%%%%%%%%%%%%%%%%%%%%%%%%%%%%%
\subsection{Likelihood-Ratio Test in Logistic Regression}
\label{subsec:lr-test-logistic}
%%%%%%%%%%%%%%%%%%%%%%%%%%%%%%%%%%%%%%%%%%%%%%%%%%%%%%%%%%%%

Suppose a reduced model is nested within a larger model.

Define

\[
\boxed{
G^2
=
2
\left[
\ell(\widehat{\boldsymbol{\beta}}_{\text{full}})
-
\ell(\widehat{\boldsymbol{\beta}}_{\text{reduced}})
\right].
}
\]

Under regularity conditions,

\[
\boxed{
G^2
\approx
\chi^2_q,
}
\]

where \(q\) is the difference in the number of free parameters.

%%%%%%%%%%%%%%%%%%%%%%%%%%%%%%%%%%%%%%%%%%%%%%%%%%%%%%%%%%%%
\subsection{Score Test in Logistic Regression}
\label{subsec:score-test-logistic}
%%%%%%%%%%%%%%%%%%%%%%%%%%%%%%%%%%%%%%%%%%%%%%%%%%%%%%%%%%%%

A score test evaluates the derivative of the log-likelihood under the
null model.

Like the Wald and likelihood-ratio tests, it is asymptotically
chi-square or normal under standard conditions.

The three procedures are asymptotically equivalent but can differ in
finite samples.

%%%%%%%%%%%%%%%%%%%%%%%%%%%%%%%%%%%%%%%%%%%%%%%%%%%%%%%%%%%%
\subsection{Logistic Regression Confidence Intervals}
\label{subsec:logistic-confidence-intervals}
%%%%%%%%%%%%%%%%%%%%%%%%%%%%%%%%%%%%%%%%%%%%%%%%%%%%%%%%%%%%

A large-sample Wald interval for \(\beta_j\) is

\[
\boxed{
\widehat{\beta}_j
\pm
z^*
SE(\widehat{\beta}_j).
}
\]

Exponentiating gives an approximate interval for the corresponding odds
ratio:

\[
\boxed{
\left(
e^{\widehat{\beta}_j-z^*SE},
\,
e^{\widehat{\beta}_j+z^*SE}
\right).
}
\]

Likelihood-profile intervals may perform better in some problems.

%%%%%%%%%%%%%%%%%%%%%%%%%%%%%%%%%%%%%%%%%%%%%%%%%%%%%%%%%%%%
\subsection{Deviance}
\label{subsec:categorical-deviance}
%%%%%%%%%%%%%%%%%%%%%%%%%%%%%%%%%%%%%%%%%%%%%%%%%%%%%%%%%%%%

For likelihood-based models, deviance compares a fitted model with a
saturated model.

A generic form is

\[
\boxed{
D
=
2
\left[
\ell_{\text{saturated}}
-
\ell_{\text{fitted}}
\right].
}
\]

Smaller deviance indicates a fitted model closer to the saturated model.

Differences in deviance are especially useful for nested-model
comparisons.

%%%%%%%%%%%%%%%%%%%%%%%%%%%%%%%%%%%%%%%%%%%%%%%%%%%%%%%%%%%%
\subsection{Saturated Model}
\label{subsec:saturated-model}
%%%%%%%%%%%%%%%%%%%%%%%%%%%%%%%%%%%%%%%%%%%%%%%%%%%%%%%%%%%%

A saturated model contains enough parameters to reproduce the observed
data exactly.

Its likelihood is therefore the largest obtainable within the response
family.

The saturated model serves as a reference model rather than usually
being a useful predictive model.

%%%%%%%%%%%%%%%%%%%%%%%%%%%%%%%%%%%%%%%%%%%%%%%%%%%%%%%%%%%%
\subsection{Null Model}
\label{subsec:null-model-logistic}
%%%%%%%%%%%%%%%%%%%%%%%%%%%%%%%%%%%%%%%%%%%%%%%%%%%%%%%%%%%%

A null or intercept-only model contains no explanatory predictors:

\[
\boxed{
\operatorname{logit}(p)
=
\beta_0.
}
\]

Comparing the fitted model with this null model measures whether the
predictors improve fit.

%%%%%%%%%%%%%%%%%%%%%%%%%%%%%%%%%%%%%%%%%%%%%%%%%%%%%%%%%%%%
\subsection{Pseudo-\(R^2\)}
\label{subsec:pseudo-r-squared}
%%%%%%%%%%%%%%%%%%%%%%%%%%%%%%%%%%%%%%%%%%%%%%%%%%%%%%%%%%%%

Several quantities called pseudo-\(R^2\) exist for logistic regression.

They are not generally equal to the ordinary least-squares

\[
R^2.
\]

Examples compare:

\[
\boxed{
\text{log-likelihoods},
}
\]

or

\[
\boxed{
\text{deviances}.
}
\]

Therefore the exact definition should always be reported.

%%%%%%%%%%%%%%%%%%%%%%%%%%%%%%%%%%%%%%%%%%%%%%%%%%%%%%%%%%%%
\subsection{Classification Thresholds}
\label{subsec:classification-threshold}
%%%%%%%%%%%%%%%%%%%%%%%%%%%%%%%%%%%%%%%%%%%%%%%%%%%%%%%%%%%%

Logistic regression estimates probabilities.

A binary classification requires an additional threshold \(c\):

\[
\widehat{Y}
=
\begin{cases}
1,&\widehat{p}\geq c,\\
0,&\widehat{p}<c.
\end{cases}
\]

The common choice

\[
c=0.5
\]

is not universally optimal.

The threshold depends on the costs of false positives and false
negatives.

%%%%%%%%%%%%%%%%%%%%%%%%%%%%%%%%%%%%%%%%%%%%%%%%%%%%%%%%%%%%
\subsection{Confusion Matrix}
\label{subsec:confusion-matrix}
%%%%%%%%%%%%%%%%%%%%%%%%%%%%%%%%%%%%%%%%%%%%%%%%%%%%%%%%%%%%

For binary classification:

\[
\begin{array}{c|cc}
&
\text{Actual }1
&
\text{Actual }0
\\
\hline
\text{Predicted }1
&
TP
&
FP
\\
\text{Predicted }0
&
FN
&
TN
\end{array}
\]

where

\[
TP
=
\text{true positives},
\]

\[
FP
=
\text{false positives},
\]

\[
FN
=
\text{false negatives},
\]

and

\[
TN
=
\text{true negatives}.
\]

%%%%%%%%%%%%%%%%%%%%%%%%%%%%%%%%%%%%%%%%%%%%%%%%%%%%%%%%%%%%
\subsection{Sensitivity and Specificity}
\label{subsec:sensitivity-specificity}
%%%%%%%%%%%%%%%%%%%%%%%%%%%%%%%%%%%%%%%%%%%%%%%%%%%%%%%%%%%%

Sensitivity is

\[
\boxed{
\operatorname{Sensitivity}
=
\frac{
TP
}{
TP+FN
}.
}
\]

Specificity is

\[
\boxed{
\operatorname{Specificity}
=
\frac{
TN
}{
TN+FP
}.
}
\]

Changing the classification threshold generally changes both.

%%%%%%%%%%%%%%%%%%%%%%%%%%%%%%%%%%%%%%%%%%%%%%%%%%%%%%%%%%%%
\subsection{Positive Predictive Value}
\label{subsec:positive-predictive-value}
%%%%%%%%%%%%%%%%%%%%%%%%%%%%%%%%%%%%%%%%%%%%%%%%%%%%%%%%%%%%

Positive predictive value is

\[
\boxed{
PPV
=
\frac{
TP
}{
TP+FP
}.
}
\]

Negative predictive value is

\[
\boxed{
NPV
=
\frac{
TN
}{
TN+FN
}.
}
\]

Unlike sensitivity and specificity under many diagnostic formulations,
predictive values depend strongly on event prevalence.

%%%%%%%%%%%%%%%%%%%%%%%%%%%%%%%%%%%%%%%%%%%%%%%%%%%%%%%%%%%%
\subsection{ROC Curve Preview}
\label{subsec:roc-curve-preview}
%%%%%%%%%%%%%%%%%%%%%%%%%%%%%%%%%%%%%%%%%%%%%%%%%%%%%%%%%%%%

A receiver operating characteristic curve plots

\[
\boxed{
\text{sensitivity}
}
\]

against

\[
\boxed{
1-\text{specificity}
}
\]

as the classification threshold varies.

It evaluates ranking or discrimination across all thresholds.

%%%%%%%%%%%%%%%%%%%%%%%%%%%%%%%%%%%%%%%%%%%%%%%%%%%%%%%%%%%%
\subsection{Area Under the ROC Curve}
\label{subsec:auc-preview}
%%%%%%%%%%%%%%%%%%%%%%%%%%%%%%%%%%%%%%%%%%%%%%%%%%%%%%%%%%%%

The area under the ROC curve, abbreviated AUC, summarizes discrimination.

One interpretation is

\[
\boxed{
AUC
=
P
(
\text{a randomly selected positive case receives a higher score than a
randomly selected negative case}
).
}
\]

An AUC near

\[
0.5
\]

corresponds to chance-level ranking under the usual orientation.

%%%%%%%%%%%%%%%%%%%%%%%%%%%%%%%%%%%%%%%%%%%%%%%%%%%%%%%%%%%%
\subsection{Calibration}
\label{subsec:logistic-calibration}
%%%%%%%%%%%%%%%%%%%%%%%%%%%%%%%%%%%%%%%%%%%%%%%%%%%%%%%%%%%%

A probability model is calibrated if predicted probabilities correspond
to observed event frequencies.

For example, among observations receiving predicted probability near

\[
0.70,
\]

approximately \(70\%\) should experience the event in a well-calibrated
model.

Discrimination and calibration are different properties.

%%%%%%%%%%%%%%%%%%%%%%%%%%%%%%%%%%%%%%%%%%%%%%%%%%%%%%%%%%%%
\subsection{Separation in Logistic Regression}
\label{subsec:logistic-separation}
%%%%%%%%%%%%%%%%%%%%%%%%%%%%%%%%%%%%%%%%%%%%%%%%%%%%%%%%%%%%

Complete separation occurs when predictor values perfectly divide the
outcome classes.

Then some maximum likelihood coefficient estimates may diverge toward

\[
\pm\infty.
\]

Ordinary asymptotic standard errors may become unreliable.

Penalized or exact procedures may be needed.

%%%%%%%%%%%%%%%%%%%%%%%%%%%%%%%%%%%%%%%%%%%%%%%%%%%%%%%%%%%%
\subsection{Overdispersion Preview}
\label{subsec:binary-overdispersion}
%%%%%%%%%%%%%%%%%%%%%%%%%%%%%%%%%%%%%%%%%%%%%%%%%%%%%%%%%%%%

For grouped binomial data, observed variability may exceed the binomial
variance

\[
np(1-p).
\]

This is called overdispersion.

Possible causes include:

\[
\boxed{
\text{unmodeled heterogeneity},
}
\]

\[
\boxed{
\text{dependence},
}
\]

or

\[
\boxed{
\text{model misspecification}.
}
\]

%%%%%%%%%%%%%%%%%%%%%%%%%%%%%%%%%%%%%%%%%%%%%%%%%%%%%%%%%%%%
\subsection{Multinomial Outcomes}
\label{subsec:multinomial-outcomes}
%%%%%%%%%%%%%%%%%%%%%%%%%%%%%%%%%%%%%%%%%%%%%%%%%%%%%%%%%%%%

Suppose the response has more than two unordered categories:

\[
Y\in\{1,\ldots,K\}.
\]

Then binary logistic regression is insufficient.

Multinomial logistic regression models probabilities

\[
P(Y=k\mid X)
\]

relative to a chosen reference category.

%%%%%%%%%%%%%%%%%%%%%%%%%%%%%%%%%%%%%%%%%%%%%%%%%%%%%%%%%%%%
\subsection{Baseline-Category Logistic Model}
\label{subsec:baseline-category-logit}
%%%%%%%%%%%%%%%%%%%%%%%%%%%%%%%%%%%%%%%%%%%%%%%%%%%%%%%%%%%%

Choose category \(K\) as the reference.

For

\[
j=1,\ldots,K-1,
\]

model

\[
\boxed{
\log
\left[
\frac{
P(Y=j\mid X)
}{
P(Y=K\mid X)
}
\right]
=
\beta_{j0}
+
\boldsymbol{\beta}_j^T X.
}
\]

Each category receives its own coefficient vector relative to the
reference category.

%%%%%%%%%%%%%%%%%%%%%%%%%%%%%%%%%%%%%%%%%%%%%%%%%%%%%%%%%%%%
\subsection{Ordinal Responses}
\label{subsec:ordinal-responses}
%%%%%%%%%%%%%%%%%%%%%%%%%%%%%%%%%%%%%%%%%%%%%%%%%%%%%%%%%%%%

If categories have a natural order, multinomial regression may ignore
useful ordering information.

Ordinal regression models cumulative probabilities while respecting the
category ordering.

A common model is the proportional-odds model.

%%%%%%%%%%%%%%%%%%%%%%%%%%%%%%%%%%%%%%%%%%%%%%%%%%%%%%%%%%%%
\subsection{Proportional-Odds Model}
\label{subsec:proportional-odds-model}
%%%%%%%%%%%%%%%%%%%%%%%%%%%%%%%%%%%%%%%%%%%%%%%%%%%%%%%%%%%%

Suppose categories are ordered

\[
1<2<\cdots<K.
\]

A cumulative logit model may be written

\[
\boxed{
\log
\left[
\frac{
P(Y\leq j\mid X)
}{
P(Y>j\mid X)
}
\right]
=
\alpha_j
-
\boldsymbol{\beta}^TX,
}
\]

for

\[
j=1,\ldots,K-1.
\]

%%%%%%%%%%%%%%%%%%%%%%%%%%%%%%%%%%%%%%%%%%%%%%%%%%%%%%%%%%%%
\subsection{Proportional-Odds Assumption}
\label{subsec:proportional-odds-assumption}
%%%%%%%%%%%%%%%%%%%%%%%%%%%%%%%%%%%%%%%%%%%%%%%%%%%%%%%%%%%%

The same coefficient vector

\[
\boldsymbol{\beta}
\]

appears for every cumulative cutpoint.

Thus predictor effects are assumed constant across the cumulative logits.

This is the proportional-odds assumption.

If it fails, more flexible ordinal models may be required.

%%%%%%%%%%%%%%%%%%%%%%%%%%%%%%%%%%%%%%%%%%%%%%%%%%%%%%%%%%%%
\subsection{Log-Linear Models}
\label{subsec:log-linear-models}
%%%%%%%%%%%%%%%%%%%%%%%%%%%%%%%%%%%%%%%%%%%%%%%%%%%%%%%%%%%%

Log-linear models analyze expected cell counts in contingency tables.

Suppose

\[
m_{ij}
=
\mathbb{E}[N_{ij}].
\]

A two-way independence model may be written

\[
\boxed{
\log m_{ij}
=
\lambda
+
\lambda_i^A
+
\lambda_j^B.
}
\]

No interaction term appears under independence.

%%%%%%%%%%%%%%%%%%%%%%%%%%%%%%%%%%%%%%%%%%%%%%%%%%%%%%%%%%%%
\subsection{Log-Linear Interaction}
\label{subsec:log-linear-interaction}
%%%%%%%%%%%%%%%%%%%%%%%%%%%%%%%%%%%%%%%%%%%%%%%%%%%%%%%%%%%%

A saturated two-way model includes

\[
\boxed{
\log m_{ij}
=
\lambda
+
\lambda_i^A
+
\lambda_j^B
+
\lambda_{ij}^{AB}.
}
\]

The interaction term

\[
\lambda_{ij}^{AB}
\]

represents association between the two categorical variables.

Thus independence corresponds to removing the interaction.

%%%%%%%%%%%%%%%%%%%%%%%%%%%%%%%%%%%%%%%%%%%%%%%%%%%%%%%%%%%%
\subsection{Connection Between Logistic and Log-Linear Models}
\label{subsec:logistic-loglinear-connection}
%%%%%%%%%%%%%%%%%%%%%%%%%%%%%%%%%%%%%%%%%%%%%%%%%%%%%%%%%%%%

Logistic regression models a conditional response probability.

Log-linear models treat all contingency-table variables symmetrically and
model expected counts.

Under appropriate sampling schemes, the two frameworks can encode the
same odds-ratio structure.

%%%%%%%%%%%%%%%%%%%%%%%%%%%%%%%%%%%%%%%%%%%%%%%%%%%%%%%%%%%%
\subsection{Stratified \(2\times2\) Tables}
\label{subsec:stratified-two-by-two}
%%%%%%%%%%%%%%%%%%%%%%%%%%%%%%%%%%%%%%%%%%%%%%%%%%%%%%%%%%%%

Suppose a possible confounding variable defines strata

\[
h=1,\ldots,H.
\]

Within each stratum, form a \(2\times2\) table:

\[
\begin{array}{c|cc}
&
Y=1
&
Y=0
\\
\hline
X=1
&
a_h
&
b_h
\\
X=0
&
c_h
&
d_h
\end{array}
\]

Association can then be examined while controlling for the stratifying
variable.

%%%%%%%%%%%%%%%%%%%%%%%%%%%%%%%%%%%%%%%%%%%%%%%%%%%%%%%%%%%%
\subsection{Mantel--Haenszel Odds Ratio}
\label{subsec:mantel-haenszel-odds-ratio}
%%%%%%%%%%%%%%%%%%%%%%%%%%%%%%%%%%%%%%%%%%%%%%%%%%%%%%%%%%%%

A common estimator of a common odds ratio across strata is

\[
\boxed{
\widehat{OR}_{MH}
=
\frac{
\sum_{h=1}^{H}
\frac{
a_hd_h
}{
n_h
}
}{
\sum_{h=1}^{H}
\frac{
b_hc_h
}{
n_h
}
}.
}
\]

This combines information across strata while accounting for the
stratification variable.

%%%%%%%%%%%%%%%%%%%%%%%%%%%%%%%%%%%%%%%%%%%%%%%%%%%%%%%%%%%%
\subsection{Mantel--Haenszel Test Preview}
\label{subsec:mantel-haenszel-test-preview}
%%%%%%%%%%%%%%%%%%%%%%%%%%%%%%%%%%%%%%%%%%%%%%%%%%%%%%%%%%%%

The Mantel--Haenszel test evaluates conditional association between two
binary variables while controlling for a categorical stratification
variable.

It is especially useful when many \(2\times2\) tables share a common
association structure.

Regression methods provide a more general framework when adjustment
involves many or continuous covariates.

%%%%%%%%%%%%%%%%%%%%%%%%%%%%%%%%%%%%%%%%%%%%%%%%%%%%%%%%%%%%
\subsection{Confounding in Categorical Data}
\label{subsec:categorical-confounding}
%%%%%%%%%%%%%%%%%%%%%%%%%%%%%%%%%%%%%%%%%%%%%%%%%%%%%%%%%%%%

A third variable \(Z\) can distort the observed association between

\[
X
\]

and

\[
Y.
\]

A crude odds ratio may differ substantially from stratum-specific or
adjusted odds ratios.

This may indicate confounding, interaction, or both.

Scientific reasoning is required to distinguish them.

%%%%%%%%%%%%%%%%%%%%%%%%%%%%%%%%%%%%%%%%%%%%%%%%%%%%%%%%%%%%
\subsection{Simpson's Paradox}
\label{subsec:simpsons-paradox-categorical}
%%%%%%%%%%%%%%%%%%%%%%%%%%%%%%%%%%%%%%%%%%%%%%%%%%%%%%%%%%%%

Simpson's paradox occurs when an association in aggregated data reverses
or changes substantially after stratification.

Symbolically, one may have

\[
\boxed{
X
\text{ positively associated with }Y
\text{ marginally}
}
\]

but

\[
\boxed{
X
\text{ negatively associated with }Y
\text{ within strata of }Z.
}
\]

This illustrates why marginal and conditional associations can differ.

%%%%%%%%%%%%%%%%%%%%%%%%%%%%%%%%%%%%%%%%%%%%%%%%%%%%%%%%%%%%
\subsection{Confounding Versus Interaction}
\label{subsec:confounding-versus-interaction-categorical}
%%%%%%%%%%%%%%%%%%%%%%%%%%%%%%%%%%%%%%%%%%%%%%%%%%%%%%%%%%%%

Confounding occurs when an additional variable distorts the relationship
of primary interest.

Interaction or effect modification occurs when the strength of
association genuinely differs across levels of another variable.

If stratum-specific odds ratios differ substantially, interaction may be
present.

If they are similar to one another but differ from the crude odds ratio,
confounding may be present.

%%%%%%%%%%%%%%%%%%%%%%%%%%%%%%%%%%%%%%%%%%%%%%%%%%%%%%%%%%%%
\subsection{Sparse Tables}
\label{subsec:sparse-contingency-tables}
%%%%%%%%%%%%%%%%%%%%%%%%%%%%%%%%%%%%%%%%%%%%%%%%%%%%%%%%%%%%

A sparse table contains many small or zero cell counts.

Asymptotic chi-square approximations may perform poorly.

Possible alternatives include:

\[
\boxed{
\text{exact tests},
}
\]

\[
\boxed{
\text{Monte Carlo tests},
}
\]

or

\[
\boxed{
\text{penalized likelihood methods}.
}
\]

%%%%%%%%%%%%%%%%%%%%%%%%%%%%%%%%%%%%%%%%%%%%%%%%%%%%%%%%%%%%
\subsection{Zero Cells}
\label{subsec:zero-cells}
%%%%%%%%%%%%%%%%%%%%%%%%%%%%%%%%%%%%%%%%%%%%%%%%%%%%%%%%%%%%

If one cell in a \(2\times2\) table is zero, the sample odds ratio

\[
\frac{ad}{bc}
\]

may become zero or infinite.

Likewise,

\[
\log\widehat{OR}
\]

may be undefined.

Small-sample corrections or exact methods can be used when appropriate.

%%%%%%%%%%%%%%%%%%%%%%%%%%%%%%%%%%%%%%%%%%%%%%%%%%%%%%%%%%%%
\subsection{Yates Continuity Correction Preview}
\label{subsec:yates-correction}
%%%%%%%%%%%%%%%%%%%%%%%%%%%%%%%%%%%%%%%%%%%%%%%%%%%%%%%%%%%%

For a \(2\times2\) Pearson chi-square test, a continuity correction may
replace

\[
|O-E|
\]

with approximately

\[
\boxed{
|O-E|-0.5.
}
\]

The correction was designed to improve approximation of discrete counts
by a continuous chi-square distribution.

It can be conservative and is not universally preferred.

%%%%%%%%%%%%%%%%%%%%%%%%%%%%%%%%%%%%%%%%%%%%%%%%%%%%%%%%%%%%
\subsection{Categorical Effect Size}
\label{subsec:categorical-effect-size}
%%%%%%%%%%%%%%%%%%%%%%%%%%%%%%%%%%%%%%%%%%%%%%%%%%%%%%%%%%%%

A chi-square test gives evidence about association but not necessarily
its magnitude.

Effect-size measures include:

\[
\boxed{
\phi,
}
\]

\[
\boxed{
\text{Cramér's }V,
}
\]

\[
\boxed{
\text{risk difference},
}
\]

\[
\boxed{
\text{relative risk},
}
\]

and

\[
\boxed{
\text{odds ratio}.
}
\]

%%%%%%%%%%%%%%%%%%%%%%%%%%%%%%%%%%%%%%%%%%%%%%%%%%%%%%%%%%%%
\subsection{Phi Coefficient}
\label{subsec:phi-coefficient}
%%%%%%%%%%%%%%%%%%%%%%%%%%%%%%%%%%%%%%%%%%%%%%%%%%%%%%%%%%%%

For a \(2\times2\) table,

\[
\boxed{
\phi
=
\sqrt{
\frac{
X^2
}{
n
}
}
}
\]

up to sign conventions when a signed binary correlation interpretation
is used.

It measures strength of association for two binary variables.

%%%%%%%%%%%%%%%%%%%%%%%%%%%%%%%%%%%%%%%%%%%%%%%%%%%%%%%%%%%%
\subsection{Cramér's \(V\)}
\label{subsec:cramers-v}
%%%%%%%%%%%%%%%%%%%%%%%%%%%%%%%%%%%%%%%%%%%%%%%%%%%%%%%%%%%%

For an \(r\times c\) table,

\[
\boxed{
V
=
\sqrt{
\frac{
X^2
}{
n
\min(r-1,c-1)
}
}.
}
\]

This generalizes the magnitude idea of \(\phi\) beyond \(2\times2\)
tables.

%%%%%%%%%%%%%%%%%%%%%%%%%%%%%%%%%%%%%%%%%%%%%%%%%%%%%%%%%%%%
\subsection{Modeling Workflow for Categorical Data}
\label{subsec:categorical-modeling-workflow}
%%%%%%%%%%%%%%%%%%%%%%%%%%%%%%%%%%%%%%%%%%%%%%%%%%%%%%%%%%%%

A useful categorical-data workflow is:

\begin{enumerate}
    \item identify whether variables are nominal, ordinal, or binary;
    \item construct frequency and contingency tables;
    \item calculate relevant marginal and conditional proportions;
    \item visualize distributions when useful;
    \item identify the scientific measure of association;
    \item decide whether inference is exact or asymptotic;
    \item inspect expected cell counts;
    \item perform an appropriate test;
    \item report effect sizes and confidence intervals;
    \item consider stratification or regression adjustment;
    \item inspect interactions;
    \item assess sparse-cell problems;
    \item interpret results according to the sampling or experimental
          design.
\end{enumerate}

%%%%%%%%%%%%%%%%%%%%%%%%%%%%%%%%%%%%%%%%%%%%%%%%%%%%%%%%%%%%
\subsection{Common Errors}
\label{subsec:categorical-common-errors}
%%%%%%%%%%%%%%%%%%%%%%%%%%%%%%%%%%%%%%%%%%%%%%%%%%%%%%%%%%%%

\subsubsection{Comparing Counts Without Considering Group Size}

Raw counts can be misleading when marginal group sizes differ.

Conditional proportions are often more meaningful.

%%%%%%%%%%%%%%%%%%%%%%%%%%%%%%%%%%%%%%%%%%%%%%%%%%%%%%%%%%%%
\subsubsection{Confusing Joint and Conditional Percentages}

The quantity

\[
P(A\cap B)
\]

differs from

\[
P(A\mid B).
\]

The denominator must always be identified.

%%%%%%%%%%%%%%%%%%%%%%%%%%%%%%%%%%%%%%%%%%%%%%%%%%%%%%%%%%%%
\subsubsection{Interpreting Association as Causation}

A contingency-table association does not automatically imply a causal
relationship.

Study design and confounding must be considered.

%%%%%%%%%%%%%%%%%%%%%%%%%%%%%%%%%%%%%%%%%%%%%%%%%%%%%%%%%%%%
\subsubsection{Using Chi-Square Approximations with Very Small Expected Counts}

Asymptotic chi-square methods may be unreliable for sparse tables.

Exact or simulation methods may be preferable.

%%%%%%%%%%%%%%%%%%%%%%%%%%%%%%%%%%%%%%%%%%%%%%%%%%%%%%%%%%%%
\subsubsection{Using Observed Counts Instead of Expected Counts to Judge the Approximation}

The quality of the chi-square approximation depends primarily on
expected counts under the null model.

%%%%%%%%%%%%%%%%%%%%%%%%%%%%%%%%%%%%%%%%%%%%%%%%%%%%%%%%%%%%
\subsubsection{Interpreting an Odds Ratio as a Relative Risk}

An odds ratio compares odds, not probabilities.

For common events,

\[
OR
\]

can differ substantially from

\[
RR.
\]

%%%%%%%%%%%%%%%%%%%%%%%%%%%%%%%%%%%%%%%%%%%%%%%%%%%%%%%%%%%%
\subsubsection{Interpreting an Odds Ratio of \(2\) as Double the Probability}

An odds ratio of \(2\) means the odds are doubled.

The corresponding probability change depends on the baseline
probability.

%%%%%%%%%%%%%%%%%%%%%%%%%%%%%%%%%%%%%%%%%%%%%%%%%%%%%%%%%%%%
\subsubsection{Ignoring Pairing}

Paired binary observations require methods such as McNemar's test rather
than an ordinary independent \(2\times2\) chi-square analysis.

%%%%%%%%%%%%%%%%%%%%%%%%%%%%%%%%%%%%%%%%%%%%%%%%%%%%%%%%%%%%
\subsubsection{Ignoring Stratification}

A crude association can hide or reverse stratum-specific relationships.

Possible confounders should be considered.

%%%%%%%%%%%%%%%%%%%%%%%%%%%%%%%%%%%%%%%%%%%%%%%%%%%%%%%%%%%%
\subsubsection{Assuming All Logistic Coefficients Are Probability Differences}

Logistic coefficients describe log-odds differences.

Exponentiated coefficients describe odds ratios.

Probability effects depend on the predictor values and baseline risk.

%%%%%%%%%%%%%%%%%%%%%%%%%%%%%%%%%%%%%%%%%%%%%%%%%%%%%%%%%%%%
\subsubsection{Using Linear Regression for Binary Outcomes Without Understanding Its Limitations}

Linear probability models can be useful in some settings, but they can
predict outside \([0,1]\) and have heteroskedastic errors.

Logistic regression is usually more natural for modeling probabilities.

%%%%%%%%%%%%%%%%%%%%%%%%%%%%%%%%%%%%%%%%%%%%%%%%%%%%%%%%%%%%
\subsubsection{Treating Pseudo-\(R^2\) as Ordinary \(R^2\)}

Different pseudo-\(R^2\) definitions measure different aspects of model
fit.

They should not be interpreted identically to OLS \(R^2\).

%%%%%%%%%%%%%%%%%%%%%%%%%%%%%%%%%%%%%%%%%%%%%%%%%%%%%%%%%%%%
\subsubsection{Ignoring Complete Separation}

Perfect classification can cause ordinary logistic MLEs to diverge.

Large coefficient estimates and standard errors may indicate separation.

%%%%%%%%%%%%%%%%%%%%%%%%%%%%%%%%%%%%%%%%%%%%%%%%%%%%%%%%%%%%
\subsubsection{Using Classification Accuracy as the Only Model Metric}

A model can achieve high accuracy when one outcome class dominates.

Probability calibration, discrimination, sensitivity, specificity, and
scientific consequences should also be considered.

%%%%%%%%%%%%%%%%%%%%%%%%%%%%%%%%%%%%%%%%%%%%%%%%%%%%%%%%%%%%
\subsection{Applications}
\label{subsec:categorical-applications}
%%%%%%%%%%%%%%%%%%%%%%%%%%%%%%%%%%%%%%%%%%%%%%%%%%%%%%%%%%%%

\begin{applicationbox}

\textbf{Medicine and Epidemiology.}
Categorical methods analyze disease status, treatment response, risk,
odds ratios, relative risks, and stratified associations.

\medskip

\textbf{Education Research.}
Researchers analyze pass/fail outcomes, program participation,
categorical survey responses, and demographic group relationships.

\medskip

\textbf{Business.}
Conversion, purchase, churn, product choice, and customer categories are
naturally modeled with categorical methods.

\medskip

\textbf{Social Science.}
Contingency tables and ordinal models analyze survey responses,
preferences, and categorical behavior.

\medskip

\textbf{Engineering.}
Defect/no-defect outcomes and categorized failure modes can be analyzed
using binomial and multinomial models.

\medskip

\textbf{Environmental Science.}
Presence/absence, habitat category, species occurrence, and thresholded
environmental outcomes often require categorical analysis.

\medskip

\textbf{Machine Learning.}
Logistic and multinomial regression are fundamental probabilistic
classification methods.

\medskip

\textbf{Experimental Design.}
Binary and categorical treatment outcomes require generalized linear
models rather than ordinary Gaussian ANOVA alone.

\end{applicationbox}

%%%%%%%%%%%%%%%%%%%%%%%%%%%%%%%%%%%%%%%%%%%%%%%%%%%%%%%%%%%%
\subsection{Exercises}
\label{subsec:categorical-exercises}
%%%%%%%%%%%%%%%%%%%%%%%%%%%%%%%%%%%%%%%%%%%%%%%%%%%%%%%%%%%%

\begin{exercise}[1]
Distinguish nominal and ordinal categorical variables.
\end{exercise}

\begin{exercise}[1]
Define a contingency table.
\end{exercise}

\begin{exercise}[1]
Define marginal and conditional proportions.
\end{exercise}

\begin{exercise}[1]
State the independence condition for two categorical variables.
\end{exercise}

\begin{exercise}[1]
State the expected-count formula under independence.
\end{exercise}

\begin{exercise}[1]
State the Pearson chi-square statistic.
\end{exercise}

\begin{exercise}[1]
Define risk difference.
\end{exercise}

\begin{exercise}[1]
Define relative risk.
\end{exercise}

\begin{exercise}[1]
Define odds.
\end{exercise}

\begin{exercise}[1]
Define an odds ratio.
\end{exercise}

\begin{exercise}[1]
Define the logit function.
\end{exercise}

\begin{exercise}[1]
State the simple logistic regression model.
\end{exercise}

\begin{exercise}[2]
A table has row total \(90\), column total \(80\), and grand total
\(240\). Compute the expected count under independence.
\end{exercise}

\begin{exercise}[2]
For the table

\[
\begin{array}{c|cc}
&
Y=1
&
Y=0
\\
\hline
X=1
&
30
&
20
\\
X=0
&
15
&
35
\end{array}
\]

compute the event risk in each group.
\end{exercise}

\begin{exercise}[2]
Compute the risk difference for the previous table.
\end{exercise}

\begin{exercise}[2]
Compute the relative risk for the previous table.
\end{exercise}

\begin{exercise}[2]
Compute the odds ratio for the previous table.
\end{exercise}

\begin{exercise}[2]
If

\[
p=0.75,
\]

compute the odds.
\end{exercise}

\begin{exercise}[2]
If the odds are \(4\), compute the corresponding probability.
\end{exercise}

\begin{exercise}[2]
If a logistic coefficient is

\[
\beta_1=0.7,
\]

write the corresponding odds ratio.
\end{exercise}

\begin{exercise}[2]
For

\[
\operatorname{logit}(p)
=
-1+0.5x,
\]

compute \(p\) when \(x=2\).
\end{exercise}

\begin{exercise}[2]
A paired binary table has

\[
b=12,
\qquad
c=4.
\]

Compute the large-sample McNemar statistic.
\end{exercise}

\begin{exercise}[2]
For an \(r\times c\) table with \(r=4\) and \(c=3\), find the
chi-square degrees of freedom.
\end{exercise}

\begin{exercise}[2]
Suppose \(X^2=18\), \(n=200\), \(r=3\), and \(c=4\). Compute
Cramér's \(V\).
\end{exercise}

\begin{exercise}[3]
Derive the expected-count formula under independence.
\end{exercise}

\begin{exercise}[3]
Show that the Pearson chi-square statistic equals the sum of squared
Pearson residuals.
\end{exercise}

\begin{exercise}[3]
Derive the \(2\times2\) odds-ratio formula

\[
\widehat{OR}
=
\frac{ad}{bc}.
\]
\end{exercise}

\begin{exercise}[3]
Show algebraically why the odds ratio approaches the relative risk when
both event probabilities are small.
\end{exercise}

\begin{exercise}[3]
Derive the approximate variance of

\[
\log\widehat{OR}.
\]
\end{exercise}

\begin{exercise}[3]
Construct a large-sample confidence interval for an odds ratio.
\end{exercise}

\begin{exercise}[3]
Derive the hypergeometric probability underlying Fisher's exact test.
\end{exercise}

\begin{exercise}[3]
Explain why McNemar's test depends only on discordant pairs.
\end{exercise}

\begin{exercise}[3]
Show that the logistic function always satisfies

\[
0<p(x)<1.
\]
\end{exercise}

\begin{exercise}[3]
Derive

\[
p
=
\frac{
e^\eta
}{
1+e^\eta
}
\]

from

\[
\log
\left(
\frac{p}{1-p}
\right)
=
\eta.
\]
\end{exercise}

\begin{exercise}[3]
Explain why

\[
e^{\beta_j}
\]

is an odds ratio in logistic regression.
\end{exercise}

\begin{exercise}[3]
Explain why a logistic regression interaction changes the odds ratio for
one predictor across values of another predictor.
\end{exercise}

\begin{exercise}[3]
Compare Wald, score, and likelihood-ratio tests in logistic regression.
\end{exercise}

\begin{exercise}[3]
Explain the meaning of deviance.
\end{exercise}

\begin{exercise}[3]
Explain why ordinary \(R^2\) does not directly extend to logistic
regression.
\end{exercise}

\begin{exercise}[3]
Explain the difference between discrimination and calibration.
\end{exercise}

\begin{exercise}[3]
Explain how complete separation affects logistic maximum likelihood.
\end{exercise}

\begin{exercise}[3]
Construct a baseline-category multinomial logit model for three outcome
categories.
\end{exercise}

\begin{exercise}[3]
Explain the proportional-odds assumption.
\end{exercise}

\begin{exercise}[3]
Show how independence in a two-way contingency table corresponds to a
log-linear model without an interaction term.
\end{exercise}

\begin{exercise}[3]
Explain how Simpson's paradox can arise from combining strata.
\end{exercise}

\begin{exercise}[3]
Distinguish confounding from effect modification.
\end{exercise}

\begin{exercise}[3]
Explain why sparse contingency tables can invalidate ordinary chi-square
approximations.
\end{exercise}

\begin{exercise}[4]
Derive the asymptotic chi-square distribution of Pearson's \(X^2\)
statistic.
\end{exercise}

\begin{exercise}[4]
Derive the likelihood-ratio \(G^2\) statistic from multinomial
likelihoods.
\end{exercise}

\begin{exercise}[4]
Study the asymptotic equivalence of Pearson \(X^2\) and likelihood-ratio
\(G^2\).
\end{exercise}

\begin{exercise}[4]
Study exact conditional inference for \(2\times2\) tables.
\end{exercise}

\begin{exercise}[4]
Investigate mid-\(p\) modifications of exact tests.
\end{exercise}

\begin{exercise}[4]
Derive the score equations for logistic regression.
\end{exercise}

\begin{exercise}[4]
Derive the Hessian and Fisher information matrix for logistic regression.
\end{exercise}

\begin{exercise}[4]
Study Newton--Raphson and iteratively reweighted least squares for fitting
logistic regression.
\end{exercise}

\begin{exercise}[4]
Investigate Firth penalized logistic regression under separation.
\end{exercise}

\begin{exercise}[4]
Study Hosmer--Lemeshow-type goodness-of-fit procedures and their
limitations.
\end{exercise}

\begin{exercise}[4]
Derive multinomial logistic regression likelihood equations.
\end{exercise}

\begin{exercise}[4]
Study adjacent-category and continuation-ratio ordinal models.
\end{exercise}

\begin{exercise}[4]
Investigate proportional-odds diagnostics.
\end{exercise}

\begin{exercise}[4]
Study hierarchical log-linear models.
\end{exercise}

\begin{exercise}[4]
Derive the Mantel--Haenszel common odds-ratio estimator.
\end{exercise}

\begin{exercise}[4]
Study the Cochran--Mantel--Haenszel test.
\end{exercise}

\begin{exercise}[4]
Investigate generalized estimating equations for correlated binary data.
\end{exercise}

\begin{exercise}[4]
Study mixed-effects logistic regression.
\end{exercise}

%%%%%%%%%%%%%%%%%%%%%%%%%%%%%%%%%%%%%%%%%%%%%%%%%%%%%%%%%%%%
\subsection{Conceptual Summary}
\label{subsec:categorical-summary}
%%%%%%%%%%%%%%%%%%%%%%%%%%%%%%%%%%%%%%%%%%%%%%%%%%%%%%%%%%%%

Categorical data are summarized through counts and proportions.

For two categorical variables, a contingency table contains cell counts

\[
\boxed{
n_{ij}.
}
\]

Under independence, expected counts are

\[
\boxed{
E_{ij}
=
\frac{
n_{i\cdot}n_{\cdot j}
}{
n
}.
}
\]

Pearson's statistic is

\[
\boxed{
X^2
=
\sum
\frac{
(O-E)^2
}{
E
}.
}
\]

For a \(2\times2\) table,

\[
\boxed{
RD
=
p_1-p_0,
}
\]

\[
\boxed{
RR
=
\frac{
p_1
}{
p_0
},
}
\]

and

\[
\boxed{
OR
=
\frac{
p_1/(1-p_1)
}{
p_0/(1-p_0)
}.
}
\]

The sample odds ratio is

\[
\boxed{
\widehat{OR}
=
\frac{ad}{bc}.
}
\]

For binary regression,

\[
\boxed{
\operatorname{logit}(p)
=
\log
\left(
\frac{
p
}{
1-p
}
\right).
}
\]

Logistic regression uses

\[
\boxed{
\operatorname{logit}(p)
=
\beta_0
+
\beta_1X_1
+
\cdots
+
\beta_pX_p.
}
\]

Exponentiating a coefficient gives an adjusted odds ratio:

\[
\boxed{
e^{\beta_j}.
}
\]

Categorical-data analysis therefore combines

\[
\boxed{
\text{probability}
+
\text{contingency tables}
+
\text{likelihood}
+
\text{hypothesis testing}
+
\text{generalized regression}.
}
\]

%%%%%%%%%%%%%%%%%%%%%%%%%%%%%%%%%%%%%%%%%%%%%%%%%%%%%%%%%%%%
\subsection{Section Connections}
\label{subsec:categorical-connections}
%%%%%%%%%%%%%%%%%%%%%%%%%%%%%%%%%%%%%%%%%%%%%%%%%%%%%%%%%%%%

\chapterconnection{
Categorical Data Analysis extends the contingency-table methods introduced
in Section~\ref{sec:hypothesis-testing} and the regression framework of
Section~\ref{sec:regression}. Binary logistic regression replaces the
Gaussian conditional-mean model with a Bernoulli probability model,
while multinomial and ordinal regression extend the approach to multiple
response categories. Log-linear models provide a complementary
likelihood-based framework for contingency-table counts. The next
section, Nonparametric Statistics, develops methods requiring fewer
distributional assumptions and introduces rank, sign, and
distribution-free procedures for situations where classical parametric
models may be unsuitable.
}

%%%%%%%%%%%%%%%%%%%%%%%%%%%%%%%%%%%%%%%%%%%%%%%%%%%%%%%%%%%%
% SECTION SOURCE TRACKING
%%%%%%%%%%%%%%%%%%%%%%%%%%%%%%%%%%%%%%%%%%%%%%%%%%%%%%%%%%%%

%------------------------------------------------------------
% 8.10 Categorical Data Analysis
%------------------------------------------------------------
%
% Source basis:
%   - Standard categorical data analysis.
%   - Standard contingency-table theory.
%   - Standard generalized linear-model theory.
%
% Used for:
%   - nominal, ordinal, and binary variables;
%   - frequency tables;
%   - contingency tables;
%   - marginal, joint, and conditional proportions;
%   - categorical independence;
%   - expected counts;
%   - Pearson chi-square tests;
%   - goodness-of-fit;
%   - likelihood-ratio chi-square statistics;
%   - Pearson residuals;
%   - 2x2 tables;
%   - risk difference;
%   - relative risk;
%   - odds and odds ratios;
%   - log odds ratios;
%   - approximate odds-ratio confidence intervals;
%   - case-control sampling;
%   - prospective and retrospective designs;
%   - Fisher exact testing;
%   - matched binary data;
%   - McNemar testing;
%   - logistic regression;
%   - logistic likelihood;
%   - odds-ratio coefficient interpretation;
%   - multiple logistic regression;
%   - categorical predictors;
%   - interactions;
%   - Wald, score, and likelihood-ratio tests;
%   - deviance;
%   - pseudo-R-squared;
%   - classification thresholds;
%   - confusion matrices;
%   - sensitivity and specificity;
%   - predictive values;
%   - ROC curves and AUC;
%   - calibration;
%   - separation;
%   - overdispersion;
%   - multinomial logistic regression;
%   - ordinal regression;
%   - proportional-odds models;
%   - log-linear models;
%   - stratified tables;
%   - Mantel--Haenszel methods;
%   - confounding and effect modification;
%   - Simpson's paradox;
%   - sparse tables;
%   - categorical effect sizes;
%   - applications and exercises.
%
% Notes:
%   - Nonparametric Statistics is developed next in Section 8.11.
%   - Multivariate Statistics later extends categorical and continuous
%     models to multiple-variable systems.
%   - Mathematical Statistics develops likelihood and asymptotic theory
%     more deeply.
%   - Statistical Learning later develops classification and predictive
%     modeling in greater detail.
%
%%%%%%%%%%%%%%%%%%%%%%%%%%%%%%%%%%%%%%%%%%%%%%%%%%%%%%%%%%%%